\documentclass[a4paper,11pt,oneside,openany]{book}

\usepackage[english]{babel}

\usepackage[utf8]{inputenc} % Consente l'uso dei caratteri accentati italiani
\usepackage[T1]{fontenc} % Consente l'uso di caratteri speciali appartenenti in genere ad alfabeti non inglesi

\usepackage{fancyhdr} % Consente il rendering della copertina

\usepackage{graphicx} % Consente l'inserimento delle figure nel documento

\usepackage{subfigure} % Consente l'inserimento di sottofigure

\usepackage{tabularx} % Consente l'uso di funzionalità aggiuntive alle tabelle
\usepackage{longtable} % Consente la scrittura delle tabelle su più pagine

\usepackage[titletoc]{appendix} % Consente di porre le appendici in indice

\usepackage{float} % Necessario per inserire immagini flottanti

\usepackage{fancyvrb} % Consente l'utilizzo di finestre Verbatim (testo cche si vuole porre così com'è, senza formattazione alcuna) all'interno di una cornice

\usepackage{amsmath} % Consente l'uso di funzioni aggiuntive per le formule
\numberwithin{equation}{section}

\usepackage{hyperref} % Indice degli argomenti e rimandi cliccabili, per una navigazione più comoda all'interno del file

\usepackage{lipsum} % Consente di inserire testo fittizio all'interno della tesi. Inutile in fase di scrittura della tesi. Eliminare questa dicitura e tutte le diciture \lorem all'interno del testo.

\usepackage{lineno}
\modulolinenumbers[1]

\usepackage[nopostdot,style=super,nonumberlist,toc, automake]{glossaries}
\makeglossaries

\setacronymstyle{long-short}
\newacronym{clup}{CLup}{Customers Line-up}
\newacronym{dpcm}{d.P.C.m}{\textit{"decreto del Presidente del Consiglio dei ministri"}}
\newacronym{gps}{GPS}{Global Positioning System}
\newacronym{fifo}{FIFO}{First In First Out}
\newacronym{api}{API}{Application Programming Interface}
\newacronym{uml}{UML}{Unified Modeling Language}
\newacronym{qos}{QoS}{Quality of Service}
\newacronym{si}{SI}{International System of Units}
\newacronym{rasd}{RASD}{Requirements Analysis and Specification Document}
 % input definitions

% Per inserire i contenuti, riferirsi ai file relativi e segnati tra parentesi graffe nei parametri \input

%*******************************************************
% Impostazioni della singola pagina e della copertina
%*******************************************************
\hypersetup{ % Visualizzazione dei collegamenti cliccabili
    colorlinks,
    citecolor=black,
    filecolor=black,
    linkcolor=black,
    urlcolor=black
}

% Imposta l'intestazione delle pagine
\fancyhf{}
\pagestyle{fancy}
\lhead{}
\chead{\leftmark}
\rhead{}
\lfoot{}
\cfoot{}
\rfoot{\thepage}

%*******************************************************
% Copertina
%*******************************************************
% Imposta la copertina
% Non modificare questa parte
\usepackage[some]{background}
\SetBgScale{1}
\SetBgColor{gray}
\SetBgAngle{0}
\SetBgOpacity{0.07}
% Non modificare questa parte

% Titolo della tesi e autore
\title{CLup – Customers Line-up \\ Requirements Analysis and Specification Document}
\author{Damiano Derin}
\begin{document}
\pagenumbering{roman}
\begin{titlepage}
    \begin{center}
    % Insert the background here (front page)
    \BgThispage
    \includegraphics[scale=0.3]{images/polimi_logo.jpg}\\
    {\LARGE {\bfseries Politecnico di Milano \\}}
    \vspace{.5cm}
    {\Large {\bfseries Department of Computer Science and Engineering \\}}
    \vspace{1.0cm}
    
    {\Large {\bfseries Software Engineering 2 \\}}
    \vspace{2.0cm}
    
    
    {\LARGE {\bfseries CLup – Customers Line-up \\
    	Requirements Analysis \\ and \\ Specification Document\\}}
    \vspace{1cm}

    {\large \today \\
    }
    % Anno Accademico 2020/2021


    % Tabella per inserire i nomi del laureando, del relatore e dell'eventuale correlatore nel modo migliore all'interno della pagina
    \vfill
    \begin{table}[h]
        {\large
            \begin{tabular}{c c c c r c c | c c l}
                & & & & & & & & & Author \\
                & & & & & & & & & \bfseries Damiano Derin \\
                & & & & & & & & & \\
                & & & & & & & & & Author \\
                & & & & & & & & & \bfseries Jas Valencic \\
            \end{tabular}
        }
    \end{table}
    \vspace{1cm}
    \end{center}
\end{titlepage}

	
\frontmatter


%*******************************************************
% Indice
%*******************************************************
\tableofcontents


%*******************************************************
% Capitoli
%*******************************************************
\mainmatter

% Non modificare questa parte
\renewcommand{\chaptermark}[1]{\markboth{\MakeUppercase{\chaptername\ \thechapter.\ #1}}{}}
% Non modificare questa parte

% Inserire tutti i capitoli necessari. Si possono inserire tutti i capitoli
% all'interno di un singolo file .tex oppure un file .tex per ogni capitolo.

\begin{linenumbers}
\chapter{Introduction}


\section{Section 01}

Section 01 example

\chapter{Conclusions}
\end{linenumbers}

%*******************************************************
% Bibliografia
%*******************************************************
\nocite{*}
\bibliography{bib/bibliography}
\bibliographystyle{plain}

%*******************************************************
% Glossario
%*******************************************************
%\glsaddall
\printglossary[type=\acronymtype,title={Glossary}]

\end{document}
